\subsection{Messverfahren}
Die Durchführung beschränkte sich auf wenige Schritte: 
\begin{itemize}
    \item die Eigenschwingungen (\(\omega_1\)) der ungekoppelten Pendel messen
    \item für drei verschiedene Kopplungsgrade waren die Eigenschwingung der symetrischen (\(\omega_1\)) und asymetrischen (\(\omega_2\)) Schwingung zu bestimmen
    \item bei Schwebungsschwingung per FFT \(\omega_1\) und \(\omega_2\) bestimmen, sowie durch eine Cursormessung im \(x(t)\) Bild die Schwebungs- und Schwingungsfrequenz
    \item Beobachtung zweier gekoppelter elektrischer Schwingkreise
\end{itemize}
Die Anwendung der FFT ist ziemlich einfach, die \verb|.tx(t)| Datei wird geladen (generational joke), daraus bestimmt der Rechner das \(\mathcal{F}[x(t)](f)=X(f)\) Bild, durch Input der Cursor-Markierungen an den Peaks wird als letzter Schritt eine Gauss-Kurve angefittet.

\subsection{Messprotokoll}
\begin{wrapfigure}{r}{0.4\textwidth}
    \includegraphics[width=7cm]{redacted.png}
    \caption{These files have been redacted}
    \label{fig:redactet}
\end{wrapfigure}
Es liegt kein dedizierte Messprotokoll vor. 
Die Auslenkungsdaten liegen in \verb|.txt| Dateien vor, außerdem existieren Bilder der FFT-Diagramme. 
\clearpage
